% !TeX spellcheck = nb_NO
% Kommentaren ovenfor gir instruks til editoren din om at den skal bruke norsk stavekontroll. Dette fungerer bl.a. i TeXstudio.

\documentclass[a4paper,norsk]{article}
% Norsk orddeling
\usepackage[norsk]{babel}
% AMS-pakkene er nyttige
\usepackage{amsmath,amsfonts,amssymb,amsthm}
% For å definere lister, som f.eks. deloppgaver
\usepackage{enumitem}
% Inneholder mye nyttig, slik som \coloneqq
\usepackage{mathtools}
% Penere kalligrafiske fonter
\usepackage[utf8]{inputenc}
\usepackage[T1]{fontenc}
\usepackage[a4paper, margin=0.5in]{geometry}
\usepackage{mathabx}
\usepackage{listings}
\usepackage{xcolor}
\usepackage{ifthen}
\usepackage{parskip}

\usepackage[most]{tcolorbox}
\tcbuselibrary{theorems, breakable, skins}
\tcbset{before upper=\par, after upper=\par} % allows itemize and enumerate inside boxes

\usepackage{hyperref}

\setlength{\parindent}{0pt}
\setlength{\parskip}{6pt plus 2pt minus 1pt}


\tcbset{
  % Better page breaking for all tcolorbox environments
  enhanced,
  breakable,
  before skip=5pt plus 2pt minus 1pt,
  after skip=5pt plus 2pt minus 1pt,
  parbox=false,
}
\newtcbtheorem[]{definitionbox}{Definisjon}{
  colback=blue!3,
  colframe=blue!70!black,
  fonttitle=\bfseries\color{white},
  coltitle=white,
  boxrule=0.6pt,
  arc=2mm,
  enhanced,
  breakable
}{def}

\newtcbtheorem[]{theorembox}{Teorem}{
  colback=green!3,
  colframe=green!50!black,
  fonttitle=\bfseries,
  coltitle=black,
  boxrule=0.6pt,
  before upper=\par\ignorespaces,
  after upper=\par\unskip,
  arc=2mm,
  enhanced,
  breakable
}{thm}

\newtcbtheorem[]{oppgavebox}{Oppgave}{
  colback=gray!5,
  colframe=black!60,
  fonttitle=\bfseries,
  coltitle=black,
  boxrule=0.6pt,
  arc=2mm,
  enhanced,
  breakable,
  parbox=false
}{oppg}

\newtcolorbox{sectionbox}[1][]{
  colback=white,
  colframe=black!70,
  boxrule=0.7pt,
  arc=3mm,
  top=2mm, bottom=2mm, left=2mm, right=2mm,
  title=#1,
  enhanced,
  breakable
}

\newtcbtheorem[]{corollarybox}{Korollar}{
  colback=purple!5,
  colframe=purple!60!black,
  fonttitle=\bfseries\color{white},
  coltitle=white,
  boxrule=0.6pt,
  arc=2mm,
  enhanced,
  breakable,
  before upper=\par\ignorespaces,
  after upper=\par\unskip
}{cor}

\newtcbtheorem[]{lemmabox}{Lemma}{
  colback=orange!5,
  colframe=orange!70!black,
  fonttitle=\bfseries,
  boxrule=0.6pt,
  arc=2mm,
  enhanced,
  breakable
}{lem}

\newtcbtheorem[]{propositionbox}{Proposisjon}{
  colback=cyan!5,
  colframe=cyan!60!black,
  fonttitle=\bfseries,
  coltitle=black,
  boxrule=0.6pt,
  arc=2mm,
  enhanced,
  breakable
}{prop}

\newtcbtheorem[]{proofbox}{Bevis}{
  enhanced,
  breakable,
  colback=white,
  colframe=black,
  fonttitle=\bfseries\color{white},
  coltitle=white,
  colbacktitle=black,
  title style={opacity=0.9},
  boxrule=0.6pt,
  arc=2mm,
  before upper=\par\ignorespaces,
  after upper=\par\unskip,
  parbox=false, % Viktig for sidebryting
  segmentation style={black,solid,opacity=0.2,line width=0.5pt}
}{proof}

\newlist{suboppgave}{enumerate}{1}
\setlist[suboppgave,1]{
  label=\textbf{(\alph*)},
  ref=\theenumi\alph*,
  leftmargin=1.5em,
  itemsep=0.4em,
  topsep=0.3em,
  parsep=0pt,
  partopsep=0pt
}



\setcounter{secnumdepth}{2}  % Sections and subsections numbered

% Python style  
\lstdefinestyle{pythonstyle}{
    language=Python,
    basicstyle=\ttfamily\small,
    keywordstyle=\color{blue},
    commentstyle=\color{green},
    numbers=left,
    frame=single
}

% Bruk bedre symboler for ≥, ≤, epsilon og phi
\renewcommand{\geq}{\geqslant}
\renewcommand{\leq}{\leqslant}
\newcommand{\eps} {\varepsilon}
\renewcommand{\epsilon}{\varepsilon}
\renewcommand{\phi}{\varphi}

% Definer de vanligste mengdene og funksjonene
\newcommand{\R}{\mathbb{R}}
\newcommand{\K}{\mathbb{K}}
\newcommand{\C}{\mathbb{C}}
\newcommand{\N}{\mathbb{N}}
\newcommand{\Z}{\mathbb{Z}}
\newcommand{\Q}{\mathbb{Q}}
% \L er definert som noe annet fra før av, så vi må redefinere \L
\renewcommand{\L}{\mathcal{L}}
% Rommet av polynomer
\newcommand{\Pol}{{\mathcal{P}}}
\DeclareMathOperator{\id}{id}
\DeclareMathOperator{\Image}{im}
\DeclareMathOperator{\Kernel}{ker}
\DeclareMathOperator{\Span}{span}
\DeclareMathOperator{\Rank}{rank}
\DeclareMathOperator{\Diag}{diag}
\DeclareMathOperator{\Proj}{proj}
\newcommand{\basis}{{\mathcal{B}}}
\newcommand{\basiss}{{\mathcal{C}}}
\newcommand{\basisss}{{\mathcal{D}}}

\newtcolorbox{algoritme}[1]{
  title=Algoritme~#1,
  colback=white,
  colframe=black,
  boxrule=0.5pt,
  arc=2mm,
  left=3mm,
  right=3mm,
  top=1mm,
  bottom=1mm
}

% Definer oppgavemiljøet
\newenvironment{oppgave}[1]%
	{\trivlist{}\item\relax\textbf{Løsning på #1.}\medspace}%
	{\endtrivlist}
% Deloppgavelister
\newlist{deloppgave}{enumerate}{1}
\setlist[deloppgave,1]{label=(\textbf{\alph*}),align=left}

\title{Real Analysis, Assignment 1 \\ MAT2400 University of Oslo, Spring 2026}
\author{Oliver Ekeberg}


\begin{document}
\maketitle
\tableofcontents


\section*{Exc 1}
\addcontentsline{toc}{section}{Exc 1}

\subsection*{1 a)}
\addcontentsline{toc}{subsection}{1 a)}


Let $ m := \inf \left\{ f(x) : x \in K \right\}  $. Argue that there exists a minimizing sequence $ \left( x_n \right)  $ in K such that $ f(x_n) \to m $ as $ n \to \infty $. 

\begin{proofbox}{1 a)}{}
    Let m exists, i.e. $ m > - \infty $. So there exists a $ x \in K $ such that $ f(x) < m + \epsilon  $. Let $ \epsilon_n := \frac{ 1 }{ n } $. This means that for an arbitrary sequence $ \left( x_n \right)  $ in K, we have that
    \[
        f(x_n)  < m + \epsilon_n = m + \frac{ 1 }{ n }.
    \]
    

    If we now let $ n \to \infty $, we get that $ f(x_n) \to m $ as $ n \to \infty $. 
    
\end{proofbox}

\subsection*{1 b)}
\addcontentsline{toc}{subsection}{1 b)}
Show that there is a subsequence $ \left( x_{n_k} \right)  $ in K, and some $ \bar{x} \in  X $ such that $ x_{n_k} \to \bar{x} $ as $ k \to \infty $.  
\begin{proofbox}{1 b)}{}
    For our arbitrary sequence $ \left( x_n \right)  $ in K, since K is compact, we now have that $ \left( x_n \right)  $ has a convergent subsequence $ \left( x_{n_k} \right)  $ converging to a point $ x \in K $. We can now just set $ x = \bar{x} $ and get that 
    \[
        x_{n_k} \to \bar{x}
    \]
    as $ k \to \infty $.
\end{proofbox}


\subsection*{1 c)}
\addcontentsline{toc}{subsection}{1 c)}

Show that $ m \leq f( \bar{x}) \leq m + \epsilon  $ for any $ \epsilon > 0 $, and conclude that $ \bar{x} $ is a minimum of f.

\begin{proofbox}{1 c)}{}
    We have that for a sequence $ \left( x_n \right)  $ in K, we have constructed a sequence $ f(x_n) $ that converges to $ m $ as $ n \to \infty $. In addition, we have that a subsequence $ \left( x_{n_k} \right)  $  converges to $ \bar{x} $ as $ k \to \infty $. 
    
    

    To prove that $ m \leq  f( \bar{x} ) \leq  m + \epsilon  $, we first show that $ m \leq f( \bar{x}) $. Since $ m := \inf \left\{ f(x) :  x \in K \right\}  $, and $ \bar(x) \in K $, then obviously, $ m \leq  f( \bar{x}) $


    To prove that $ f( \bar{x}) \leq  m+ \epsilon  $, we use lower semicontinuity. Since $ x_{n_k} \to \bar{x} $ as $ k \to \infty $, we can for every $ \delta > 0 $  pick a $ k \geq  N \in \mathbb{N} $ such that $ d \left( \bar{x}, x_{n_k} \right) < \delta \Rightarrow f( \bar{x}) < f(x_{n_k}) + \epsilon   $. Since we allow $ \left( x_{n_k} \right)  $ to be a minimizing sequence, we get that 
    \[
        f(\bar{x}) < \lim_{k \to \infty} f( x_{n_k}) + \epsilon _k =  \lim_{k \to \infty} f( x_{n_k}) + \frac{ 1 }{ k }.
    \]

    This means that $ f (\bar{x}) \leq  m $. Hence, $ m \leq  f (\bar{x}) \leq m $  and we conclude that $ \bar{x} $ is a minimizer of f.
    
\end{proofbox}


\begin{proofbox}{1 c)}{}
  Show that $ m \leq  f(\bar{x} \leq  m + \epsilon  $. Since $ \bar{x} \in K $, we have that obviously, the infimum of a set, is smaller than or equal to any possible value in that set. Hence,
  \[
    m \leq f(\bar{x}.
  \]
  

  To show that $ f \left( \bar{x} \right) \leq  m + \epsilon   $, we let $ \left( x_n \right)  $ be a minimizing sequence of f. Since $ \left( x_n \right)  $ is a sequence in a compact metric space, we have that it has a convergent subsequence $ \left( x_{n_k} \right)  $ converging to $ \bar{x} \in K $. Since the initial sequence $ \left( x_n \right)  $ is a minimizing sequence, then $ \left( x_{n_k} \right)  $ is also a minimizing sequence, which means that
  \[
      f(\bar{x}) \leq \liminf_{n \to \infty} f \left( x_{n_k} \right).
  \]
  

  But since any subsequence converges to what its oversequence converges to, and $ \left( x_n \right)  $ converges to m, then 
  \[
      f(\bar{x_n}) \leq \liminf_{n \to \infty} f \left( x_{n_k} \right) \to m.
  \]
  Hence we have also proved that $ f(\bar{x} \leq  m $. This means that $ m \leq  f(\bar{x}) \leq  m $, and we can conclude that m is the minimum of f, and this is attained at $ \bar{x} $ 
\end{proofbox}

\section*{Exc 2}
\addcontentsline{toc}{section}{Exc 2}
Let $ \left( X, d \right)  $ be a metric space and $ E \subseteq X $ any subset. Let $ f : E \to \mathbb{R} $ be any bounded function (i.e., there is some $ M \in \mathbb{R} $ such that $ \left| f(x) \right| \leq M \; \forall x \in E$  ). Define $ g : \bar{E} \to \mathbb{R} $ by
\[
    g(x) := \lim_{y \to x, y \in E} \inf f(y) \; \forall x \in \bar{E}
\]

that is
$ g(x) := \lim_{n \to \infty} f_n(x)  $, where $ f_n(x) := \inf \left\{ f(y) : y \in  E \cap B \left( x; \frac{ 1 }{ n } \right)  \right\}  $

Here $ \bar{E} = E \cup \partial E $ is the closure of E. g is called the lower semi continuous envelope of f.

\subsection*{2 a)}
\addcontentsline{toc}{subsection}{2 a)}

Show that g is well defined, i.e. the above limit exists, and that $ g(x) \leq f(x) $ for all $ x \in E $.


\begin{proofbox}{2 a)}{}
    g is well defined if
    \begin{enumerate}
        \item for all $ x \in \bar{E} $, the value $ g(x) $ actually makes sense
        \item The value $ g(x) $ is an element in $ \mathbb{R} $ 
        \item $ g(x) $ is unique in $ \mathbb{R} $.
    \end{enumerate}


    Lets first establish what we are dealing with. We have a function $ g : \bar{E} \to \mathbb{R} $ defined by 
    \[
        g(x) := \lim_{n \to \infty} f_n(x)  
    \]

    where $ f_n(x) := \inf \left\{ f(y) : y \in E \cap B \left( x ; \frac{ 1 }{ n } \right)  \right\}  $. Assume that $ \left( x_m \right)  $ is a sequence that converges to a point $ x \in \bar{E} $. Then around every single $ x_m $, we can define a ball $ B \left( x_m ; \frac{ 1 }{ n } \right)  $. Then $ E \cap B \left( x_m ; \frac{ 1 }{ n } \right) \neq \emptyset  $, which means that if $ y \in E \cap B \left( x_m ; \frac{ 1 }{ n } \right) $, then $ f_n(x_m) < \infty $, because infimum of the empty set is $ \infty $. If we define
    \[
        E_n(x_m) = E \cap B \left( x_m; \frac{ 1 }{ n } \right) 
    \]
    as a sequence of sets, then the \textit{size} of the sets are decreasing. If we now define 

    \[
        f_n(x_m) =  \inf(E_n(x_m))
    \]
    then the for each n, we can define $ m := \inf(E_n(x_m)) $, where m are real numbers. This means that $ g(x) = \lim_{n \to \infty}  f_n(x) $ is a sequence of real numbers. This is well defined.
    
    

    To now show that $ g(x) \leq f(x) \; \forall x \in E $, observe that we are taking infimum over smaller and smaller sets. This means that $ f_n(x)  $ is an increasing sequence of infimum of sets. To illustrate, consider $ n \in \mathbb{N} $ and $ n > 0 $. Then $ A_n \left[ \frac{ -1 }{ n }, \frac{ 1 }{ n } \right]  $ is a sequence of sets. We can take the infimum of this set, and get that 
    \[
        \inf A_n = \frac{ -1 }{ n }.
    \]
    So the smaller larger n is, the smaller the infimum gets, i.e. $ A_{n+1} \subseteq A_n$ for all $ n \in \mathbb{N} $. 

    
    Going back to our exercise, we now see that 
    \[
        f_{n}(x) \leq  f_{n-1}(x) \leq ... \leq f(x) \; \forall n \in \mathbb{N}.
    \]
    Hence, since $ g(x) := \lim_{n \to \infty} f_n(x) \leq \lim_{n \to \infty} f(x) = f(x)   $. So $ g(x) \leq  f(x) $. We get that the limit is well defined, because the sequence $ f_n(x) $ is a monotone increasing sequence that is bounded. Hence, it is convergent.
\end{proofbox}


\subsection*{2 b)}
\addcontentsline{toc}{subsection}{2 b)}

Show that g is lower semicontinuous

\begin{proofbox}{2 b)}{}
  We can start by observing that $ f_n(x) \leq g(y)  $ for all $ y \in  E \cap B \left( x; \frac{ 1 }{ n } \right)  $. If we let $ y \in E \cap B \left( x; \frac{ 1 }{ n } \right)  $, then we have that for large enough k

\[
    E \cap B \left( y ; \frac{ 1 }{ k } \right) \subseteq E \cap B \left( x; \frac{ 1 }{ n } \right)
\]

So $ f_n(x) \leq f_k(y) $. If we now let $ k \to \infty $, we get that 
\[
    f_n(x) \leq g(y) := \lim_{k \to \infty}  f_k(y) \; \forall  y \in E \cap B \left( x; \frac{ 1 }{ n } \right) 
\]


Now observe since $ g(x) := \lim_{n \to \infty} f_n(x)  $, we have that for all $ n \geq N \in \mathbb{N} $ and $ \epsilon > 0 $ that $ \left| g(x) - f_n(x) \right| < \epsilon  $. This means that for a general $ x \in \bar{E} $ and all $ n \geq N \in \mathbb{N} $ , we have that
\[
    \left| g(x) - f_n(x) + f_n(x) \right| \leq \left| g(x) - f_n(x)  \right| + \left| f_n(x) \right| < \epsilon + f_n(x) \leq g(y) + \epsilon.
\]
Hence, g is lower semi continuous.
\end{proofbox}

\subsection*{2 c)}
\addcontentsline{toc}{subsection}{2 c)}

Show that if f is lower semicontinuous on E, then $ g(x) = f(x)  $ for all $ x \in E $.

\begin{proofbox}{2 c)}{}
  Let $ x \in E $. Since f is lower semi continuous, then $ \forall \epsilon > 0 $ and $ \delta > 0 $, $ d \left( x, y \right) < \delta  $ implies that $ f(x) < f(y) + \epsilon  $. In other words, we have that 
  \[
      f(x) \leq  liminf_{y \to x} f(y)
  \]

  Now fix $ \epsilon > 0 $, by lower semi continuity, pick $ \delta > 0 $ such that $ y \in E \cap B \left( x; \delta  \right)  $. This implies that $ f(x) < f(y) + \epsilon  $. let $ \frac{ 1 }{ N } < \delta  $, then for all $ n \geq N $, we have for every $ y \in E \cap B \left( x; \frac{ 1 }{ n } \right)  $ satisfying $ f(x) < f(y) + \epsilon  $. Since $ f_n(x) $ is an increasing sequence, we have that $ f(x) - \epsilon  \leq f_n(x) $. Taking $ n \to \infty $, we get that 
  \[
      f(x) - \epsilon  \leq g(x).
  \]
  And since $ \epsilon >0 $ is arbitrarily chosen, we get that 
  \[
      f(x) \geq  g(x).
  \]
  
  Hence, $ g(x) \geq f(x)  $ and $ f(x) \geq g(x) $, then $ g(x) = f(x) $ is the only possibility
\end{proofbox}

\end{document}